\section{Ordered elliptic sources and fixed-distance fibres}
\label{sec:elliptic-sources}

We label the vertices $00,10,01,11$ by $0,1,2,3$, respectively.
The distance surface is given by
\begin{equation}\label{eq:ell-original}
 r_0^2=x^2+y^2,\quad r_1^2=(x-1)^2+y^2,\quad
 r_2^2=x^2+(y-1)^2,\quad r_3^2=(x-1)^2+(y-1)^2.
\end{equation}
Throughout, a strict rational point has $0<x,y<1$ and
$r_0,r_1,r_2,r_3\in\Q_{>0}$.  This section gives two descriptions of
the complete source of such points: an ordered pair on a congruent-number
elliptic curve, with one final square condition, and a genus-five curve
obtained by fixing a signed pair of distances.

\subsection{Diagonal sums and a simultaneous square model}

Set
\[
 s=r_0+r_3,\qquad d=r_0-r_3,\qquad
 \tau=r_1+r_2,\qquad e=r_2-r_1.
\]
Subtracting the equations in \eqref{eq:ell-original} gives
\[
 sd=2(x+y-1),\qquad \tau e=2(x-y).
\]
The four equations are equivalently encoded by
\begin{equation}\label{eq:ell-diagonal}
 s^2+d^2=\tau^2+e^2,\qquad
 (s^2-2)(2-d^2)=\tau^2e^2,
\end{equation}
together with the return formulas
\begin{equation}\label{eq:ell-diagonal-return}
 \begin{aligned}
 (r_0,r_1,r_2,r_3)&=\tfrac12(s+d,\tau-e,\tau+e,s-d),\\
 x&=\tfrac12+\tfrac14(sd+\tau e),\qquad
 y=\tfrac12+\tfrac14(sd-\tau e).
 \end{aligned}
\end{equation}
For example, eight times the difference between the first two sides of
$x^2+y^2=r_0^2$ is
$-(s^2-2)(2-d^2)+\tau^2e^2$; the other three equations follow by
subtraction and the first equation of \eqref{eq:ell-diagonal}.

\begin{lemma}\label{lem:ell-diagonal-exclusion}
Neither diagonal of the square contains a strict rational four-distance
point.  The square symmetries let us choose the chamber
\begin{equation}\label{eq:ell-chamber}
 x+y>1,\qquad x>y.
\end{equation}
\end{lemma}
\begin{proof}
On $x=y$
one has $r_0=\sqrt2x$, and on $x+y=1$ one has
$r_1=\sqrt2(1-x)$.  The square symmetries act on
$(x+y-1,x-y)$ by signed permutations.
\end{proof}
In \eqref{eq:ell-chamber} we have
\[
 s>d>0,\qquad \tau>e>0,\qquad s,\tau>\sqrt2.
\]
The strict triangle inequalities give $d^2,e^2<2$.  In fact
\begin{equation}\label{eq:ell-J-negative}
 d^2+e^2\leq
 2\bigl((x+y-1)^2+(x-y)^2\bigr)
 =4\bigl((x-\tfrac12)^2+(y-\tfrac12)^2\bigr)<2.
\end{equation}

Define
\begin{equation}\label{eq:ell-pure-coordinates}
 p=\frac{2-e^2}{d^2},\quad
 \beta=\frac{s^2}{2-e^2},\quad
 U=\frac{s}{d},\quad V=\frac{2e}{d},\quad
 W=\frac{2\tau}{d},\quad T=\frac2d.
\end{equation}
Here $p,\beta>1$, and \eqref{eq:ell-diagonal} becomes
\begin{equation}\label{eq:ell-pure}
 \begin{aligned}
 U^2&=p\beta,\qquad
 V^2=2(p-1)(\beta-1),\\
 W^2&=2(p+1)(\beta+1),\qquad
 T^2=(p-1)\beta+p+1.
 \end{aligned}
\end{equation}
Conversely these four equations give
\begin{equation}\label{eq:ell-pure-return}
 \begin{aligned}
 (r_0,r_1,r_2,r_3)
 &=\left(\frac{U+1}{T},\frac{W-V}{2T},
          \frac{W+V}{2T},\frac{U-1}{T}\right),\\
 x&=\frac12+\frac{U+WV/4}{T^2},\qquad
 y=\frac12+\frac{U-WV/4}{T^2}.
 \end{aligned}
\end{equation}
Thus all four square roots remain attached to the same $p,\beta$.

\subsection{The complete ordered elliptic source}

For a positive squarefree integer $n$, let
\[
 E_n:\quad Y^2=X(X-n)(X+n).
\]
An ordered source consists of two points
$Q=(a,b)$ and $R=(c,\ell)$ on $E_n$, and three rational roots
$u,v,h$ satisfying
\begin{equation}\label{eq:ell-ordered-roots}
 u^2=-\frac{ac}{n},\qquad
 v^2=-\frac{(a-n)(c-n)}{n},\qquad
 h^2=\frac{a+c}{2(c+n)}.
\end{equation}
The third product root is fixed by the ordinates:
\begin{equation}\label{eq:ell-w-ordinate}
 w=\frac{b\ell}{nuv},\qquad w^2=(a+n)(c+n).
\end{equation}
In particular the first two equations specify the ordered square classes
\[
 [ac]=[-n],\qquad [(a-n)(c-n)]=[-n]
 \quad\text{in }\Q^*/\Q^{*2}.
\]
The condition on $h$ is a further square condition on this same pair.

\begin{theorem}[Exact ordered-source correspondence]
\label{thm:ell-exact-source}
Strict rational points in the chamber \eqref{eq:ell-chamber} are in
bijection with the union, over positive squarefree $n$, of the rational
solutions of \eqref{eq:ell-ordered-roots} for which
\begin{equation}\label{eq:ell-source-signs}
 -n<a<0<c-n,\qquad b,\ell,u,v,h>0,
\end{equation}
and
\begin{equation}\label{eq:ell-strict-metric}
 \frac{2u}{a+c}\left(-\frac ba+\frac\ell c\right)<2.
\end{equation}
For this correspondence the final square is
\begin{equation}\label{eq:ell-final-square}
 f(Q,R):=\frac{a+c}{2(c+n)}
       =\frac{1}{(r_0+r_3)^2}.
\end{equation}
\end{theorem}

\begin{proof}
We first give rational maps that preserve the actual roots and ordinates.
From an ordered source set
\begin{equation}\label{eq:ell-forward-pair}
 \begin{aligned}
 p&=-\frac na,& k&=\frac b{a^2},&
 \beta&=\frac{c+n}{c-n},\\
 U&=\frac{(n-a)w}{bv},&
 V&=-\frac{2b}{av},&
 W&=\frac{2(a-n)u}{av},& T&=2Uh.
 \end{aligned}
\end{equation}
The curve equations and \eqref{eq:ell-ordered-roots} imply
\eqref{eq:ell-pure} and
\begin{equation}\label{eq:ell-k}
 nk^2=p(p^2-1).
\end{equation}
For instance, $k^2=b^2/a^4=(a^2-n^2)/a^3$ gives
\eqref{eq:ell-k}; the other identities follow by cancelling the factors
$a,a\pm n,c,c\pm n$.
Conversely, from \eqref{eq:ell-pure} and \eqref{eq:ell-k} put
\begin{equation}\label{eq:ell-inverse-pair}
 \begin{aligned}
 a&=-\frac np,& b&=ka^2,&
 c&=n\frac{\beta+1}{\beta-1},&
 \ell&=\frac nk\frac{UVW}{(\beta-1)^2},\\
 u&=\frac{(p-1)W}{kV},&
 v&=\frac{2(p^2-1)}{kV},&
 h&=\frac T{2U}.&&
 \end{aligned}
\end{equation}
These formulas give, separately,
\[
 \begin{aligned}
 b^2&=\frac{n^3(p^2-1)}{p^3}=a(a-n)(a+n),\\
 \ell^2&=\frac{4n^3\beta(\beta+1)}{(\beta-1)^3}
          =c(c-n)(c+n),\\
 u^2&=\frac{n(\beta+1)}{p(\beta-1)}=-\frac{ac}{n},\\
 v^2&=\frac{2n(p+1)}{p(\beta-1)}=-\frac{(a-n)(c-n)}n,\\
 h^2&=\frac{(p-1)\beta+p+1}{4p\beta}
       =\frac{a+c}{2(c+n)}.
 \end{aligned}
\]
Inserting \eqref{eq:ell-inverse-pair} into
\eqref{eq:ell-forward-pair} returns $U,V,W,T$ with their signs.
Conversely it returns $b$ and $\ell$ as well as $a,c,u,v,h$.

To obtain $n$ and $k$ from a strict original point, set
\[
 t=r_0r_3+r_1r_2=\frac{s^2-e^2}{2},\quad
 J=d^2+e^2-2,\quad K=2-d^2,
\]
and choose the unique positive squarefree integer $n$ and positive
$q\in\Q$ with $t^2-1=nq^2$.  This is possible because
\eqref{eq:ell-diagonal} gives
\begin{equation}\label{eq:ell-theta-id}
 t^2-1=\left(\frac{de}{J}\right)^2K(2-e^2)>0.
\end{equation}
For completeness, writing $D=d^2$ and $E=e^2$, the two equations of
\eqref{eq:ell-diagonal} give
\[
 s^2=-\frac{(2-E)(2+E-D)}{D+E-2},\qquad
 \tau^2=-\frac{(2-D)(2+D-E)}{D+E-2}.
\]
These expressions prove \eqref{eq:ell-theta-id} and the identity
\[
 \frac{p(p^2-1)}{t^2-1}
   =\left(\frac{\tau J^2}{d^4eK}\right)^2.
\]
Consequently
\[
 k=\frac{q\tau J^2}{d^4eK}>0
\]
satisfies \eqref{eq:ell-k}.  Equations
\eqref{eq:ell-pure-coordinates} and \eqref{eq:ell-inverse-pair}
now give the required source.  All their denominators are nonzero
by \eqref{eq:ell-J-negative}, $d,e>0$, and $p,\beta>1$.

On the source, direct use of \eqref{eq:ell-forward-pair} gives
\begin{equation}\label{eq:ell-metric-identities}
 q=\frac{2u}{a+c},\qquad
 t=\frac{c-a}{c+a},\qquad
 \frac{WV}{T^2}=-\frac{qb}{a},\qquad
 \frac{4U}{T^2}=\frac{q\ell}{c}.
\end{equation}
Thus \eqref{eq:ell-strict-metric} is precisely
$(4U+WV)/T^2<2$.  The signs in \eqref{eq:ell-source-signs}
give $U>1$, $W>V>0$ and $T>0$.  All four radii in
\eqref{eq:ell-pure-return} are positive.  That formula shows that
the metric inequality is equivalent to $x<1$ in the chamber;
it also gives $x>1/2$ and $|y-1/2|<1/2$.
Finally $r_0+r_3=2U/T=1/h$ proves
\eqref{eq:ell-final-square}.  The displayed inverse formulas and the
uniqueness of $n,q>0$ prove bijectivity.
\end{proof}

The formulas above also describe the signed algebraic covers.  Let $S$
denote the surface \eqref{eq:ell-original}, and let $X_n$ be the
normalization of $E_n\times E_n$ in the field obtained by adjoining
$u,v,h$ from \eqref{eq:ell-ordered-roots}.
\begin{proposition}\label{prop:ell-original-cover}
There is an equality of function fields
\[
 \Q(X_n)=\Q(S)(q),\qquad q^2=\frac{t^2-1}{n},
 \qquad t=r_0r_3+r_1r_2.
\]
The extension has degree two.
\end{proposition}
\begin{proof}
The four radicands in \eqref{eq:ell-original} have independent
square classes over $\overline\Q(x,y)$: each factors into two lines,
and each of these eight lines occurs for only one corner.  Hence
$S$ is geometrically integral.
The preceding maps are rational identities before signs are chosen.
Their inverse recovers $q$ from the simultaneous-square coordinates by
\[
 q=\frac{kVW}{(p^2-1)T^2}.
\]
The two root classes defining $u,v$ are geometrically independent:
their half-divisor classes on $E_n\times E_n$ restrict on either
factor to the independent classes
$[(0,0)-O]$ and $[(n,0)-O]$ in $E_n[2]$.
Their degree-four cover is therefore connected and unramified.
The zero divisor $a+c=0$ has simple generic multiplicity, which
remains odd after this unramified pullback.  Adjoining $h$ is
therefore ramified there, so this last quadratic extension is also
connected.  These statements hold over $\overline\Q$.
The resulting involution
\[
 (a,b,c,\ell,u,v,h)\longmapsto(a,-b,c,-\ell,-u,-v,h)
\]
fixes $U,V,W,T$, hence fixes $S$, and sends $q$ to $-q$.
It is nontrivial, proving the degree assertion.  The normalization of
the graph gives the corresponding birational identification of proper
models.
\end{proof}

\subsection{Fixing an opposite pair of distances}

Let $c_0=\mu/\nu$ be a fixed positive rational number, with coprime
positive integers $\mu,\nu$.  Suppose
\[
 c_0=r_3+\varepsilon r_0,\qquad
 d_0=r_3-\varepsilon r_0,\qquad \varepsilon\in\{1,-1\};
\]
for a difference we orient the pair so that $r_3>r_0$.
The strict real bounds give $0<c_0<2$.
With
\[
 A=2\nu^2-\mu^2,\quad B=4\nu^2-\mu^2,\quad
 P=\nu d_0,\quad Z=\nu^2\bigl((1-2x)-(1-2y)\bigr),
\]
and $Q=2\nu r_1$, $R=2\nu r_2$, the full equations are
\begin{equation}\label{eq:ell-fixed-opposite}
 Z^2=A(P^2-2\nu^2),\qquad
 Q^2=P^2+\mu^2+2Z,\qquad R^2=P^2+\mu^2-2Z.
\end{equation}
Indeed, if $u_0=1-2x$ and $v_0=1-2y$, then
\[
 c_0d_0=u_0+v_0,\qquad
 c_0^2+d_0^2=2+u_0^2+v_0^2,
\]
so $(u_0-v_0)^2=(2-c_0^2)(d_0^2-2)$.
The inverse is
\begin{equation}\label{eq:ell-fixed-opposite-return}
 \begin{aligned}
 x&=\frac12-\frac{\mu P+Z}{4\nu^2},&
 y&=\frac12-\frac{\mu P-Z}{4\nu^2},\\
 r_3&=\frac{\mu+P}{2\nu},&
 r_0&=\frac{\varepsilon(\mu-P)}{2\nu},&
 r_1&=\frac{|Q|}{2\nu},&r_2&=\frac{|R|}{2\nu}.
 \end{aligned}
\end{equation}
Its strict locus is specified by positivity of these radii and
\begin{equation}\label{eq:ell-fixed-strict}
 |\mu P+Z|<2\nu^2,\qquad |\mu P-Z|<2\nu^2.
\end{equation}
Substitution into \eqref{eq:ell-original} proves the inverse.

\begin{proposition}\label{prop:ell-opposite-genus}
For every rational $0<c_0<2$, the smooth projective curve determined
by \eqref{eq:ell-fixed-opposite} is geometrically connected of genus five.
\end{proposition}
\begin{proof}
Neither $A$ nor $B$ is zero.  Over the $Z$-line the three radicands,
after multiplication by the same nonzero constant $A$, are
\[
 f_0=Z^2+2A\nu^2,\quad
 f_+=(Z+A)^2+2A\mu^2,\quad
 f_-=(Z-A)^2+2A\mu^2.
\]
All three have two simple roots.  A common root of $f_0,f_+$
would be $Z=-\mu^2/2$, where $f_0=B^2/4\ne0$;
the pair $f_0,f_-$ is identical after $Z\mapsto-Z$.
A common root of $f_+,f_-$ would be $Z=0$, where
$f_+=A(2\nu^2+\mu^2)\ne0$.
Thus there are six distinct branch values.  The three square classes
are independent over $\overline\Q(Z)$, since each has simple zeros
belonging to no other radicand.  The cover has degree eight, with
four points of ramification index two above each branch value;
infinity is unramified.  Riemann--Hurwitz gives
$2g-2=-16+6\cdot4=8$.
\end{proof}

Multiplication of the last two equations in
\eqref{eq:ell-fixed-opposite} produces an elliptic quotient:
\begin{equation}\label{eq:ell-fixed-elliptic}
 F_{\mu,\nu}:\quad
 \eta^2=\xi^3+\alpha\xi^2+B^2\xi,
 \qquad \alpha=2(3\mu^2-4\nu^2),
 \qquad (\xi,\eta)=(P^2,PQR).
\end{equation}
Its cubic discriminant is nonzero because
$\alpha^2-4B^2=-32\mu^2A\ne0$.
Every strict image satisfies
\begin{equation}\label{eq:ell-quotient-interval}
 0<\xi<B.
\end{equation}
For the upper bound use $c_0^2+d_0^2<4$.
For the lower bound, $P=0$ for a sum would place the original point
on $x+y=1$, while for a difference $P=\nu(r_3+r_0)>0$.
The quotient always has the four torsion points
\begin{equation}\label{eq:ell-order-four}
 O,\quad(0,0),\quad(B,2\mu B),\quad(B,-2\mu B).
\end{equation}
The duplication formula
\begin{equation}\label{eq:ell-duplication}
 \xi(2S)=\frac{(\xi(S)^2-B^2)^2}{4\eta(S)^2}
\end{equation}
shows that the last two double to $(0,0)$.

\begin{theorem}[Uniform torsion for odd pair parameters]
\label{thm:ell-torsion}
If $\mu,\nu$ are odd, the rational torsion subgroup of
$F_{\mu,\nu}$ is cyclic of order four, with points
\eqref{eq:ell-order-four}.  Any strict four-distance point with
opposite-pair value $\mu/\nu$ gives a point of infinite order on
$F_{\mu,\nu}$.
\end{theorem}
\begin{proof}
Put $S=(0,0)$ and $T=(B,2\mu B)$. The quadratic factor
$\xi^2+\alpha\xi+B^2$ has discriminant
$32\mu^2(\mu^2-2\nu^2)$, of exact $2$-adic valuation five.
Consequently $S$ is the unique nonzero rational point killed by two.
The duplication formula gives $2T=S$, so $T$ has exact order four.
We have
\[
 B\equiv3\pmod8,\qquad \alpha\equiv14\pmod{16},\qquad
 B^2\equiv9\pmod{16}.
\]

We first exclude odd torsion. For the integral Weierstrass equation over
$\Z_2$, let $G_1$ be the subgroup of $F_{\mu,\nu}(\Q_2)$ reducing
to infinity. Its parameter $z=-\xi/\eta$ lies in $2\Z_2$.
The formal group law has integral coefficients and linear part
$z_1+z_2$. Indeed, in coordinates $z=-\xi/\eta$ and $w=-1/\eta$
the equation has a unit derivative in $w$ at $(0,0)$; successive
coefficient comparison solves for $w$ as an integral power series in
$z$. Substitution into the addition law gives
\[
 [m]z=mz+z^2g_m(z),\qquad g_m(z)\in\Z_2[[z]].
\]
If $m$ is odd and $z\ne0$, the two terms on the right have different
valuations, with the first strictly smaller. Thus $G_1$ has no
nonzero odd torsion. This argument uses the smooth point at infinity
and applies to the present integral equation at $2$.

For an arbitrary rational point $P$, put $Q=2P$. If $Q$ is finite,
its abscissa $\xi$ is a rational square by
\eqref{eq:ell-duplication}. If $\xi=0$, then $Q=S$ and $4P=O$.
Suppose $\xi\ne0$, and write $v=\ord_2(\xi)$.
The case $v=0$ is impossible. The curve equation makes the ordinate
integral, and a unit square is $1$ or $9$ modulo $16$, so
\[
 \xi(\xi^2+\alpha\xi+B^2)\equiv8\pmod{16},
\]
whereas a square in $\Z_2$ cannot have this residue.

If $v>0$, the term $B^2\xi$ has uniquely least valuation in the
curve equation. Hence $2\ord_2(\eta)=v$, and
\eqref{eq:ell-duplication} gives
$\ord_2(\xi(2Q))=-v-2\le-4$.
If $v<0$, the cubic term has uniquely least valuation, so
$2\ord_2(\eta)=3v$ and
$\ord_2(\xi(2Q))=v-2\le-4$.
The numerator and denominator of the duplication formula are nonzero
in both cases. A point with a pole in its abscissa belongs to $G_1$.
We have therefore proved that $4P\in G_1$ for every rational point
$P$, including the cases $4P=O$.
If $P$ has odd order, then $4P=O$ by the preceding formal-group
argument, and coprimality with four gives $P=O$.

Every rational torsion point has $2$-power order. If $4P=O$, then
$2P$ is $O$ or $S$. In the first case $P$ is $O$ or $S$; in the
second, $2(P-T)=O$, so $P$ is $T$ or $-T$.
There is no point of order eight: a half of $T$ or $-T$ would make
$B$ a rational square by \eqref{eq:ell-duplication}, contrary to
$B\equiv3\pmod8$. A point of larger $2$-power order would have a
multiple of order eight. Thus the torsion subgroup is precisely
$\{O,S,T,-T\}$.
Finally, every strict image satisfies
\eqref{eq:ell-quotient-interval}, whereas the finite torsion abscissas
are $0$ and $B$. This proves the infinite-order assertion.
\end{proof}

\begin{corollary}\label{cor:ell-rank-zero}
If $\mu,\nu$ are odd and $\rank F_{\mu,\nu}(\Q)=0$, then
$\mu/\nu$ is neither an opposite-pair sum nor a nonzero absolute
difference at a strict rational four-distance point.
\end{corollary}
\begin{proof}
Such a point would give a rational point of $F_{\mu,\nu}$ in the interval
\eqref{eq:ell-quotient-interval}. Rank zero makes every rational point
torsion, whereas Theorem~\ref{thm:ell-torsion} places all torsion points
outside that interval.
\end{proof}

\subsection{Adjacent pairs and fixed-pair finiteness}

For an adjacent pair choose signed radii $R_0,R_1$ with
$R_0+R_1=c_0>0$, where $R_i^2=r_i^2$.  Write $\lambda=c_0^2$.
A strict adjacent sum is greater than one and an absolute difference
is less than one, so a nonzero permitted value has
$\lambda\ne0,1$.
With $U_0=2x-1$ and $V_0=2y$, the first two distance equations
give the conic
\begin{equation}\label{eq:ell-adjacent-conic}
 (\lambda-1)U_0^2+\lambda V_0^2
       -\lambda(\lambda-1)H_0^2=0.
\end{equation}
It is smooth and contains $(c_0,0,1)$.  In the chart $H_0=1$,
\[
 R_0=\frac{c_0+U_0/c_0}{2},\qquad
 R_1=\frac{c_0-U_0/c_0}{2}.
\]
The remaining roots are given by
\[
 f_2=\frac{(U_0+H_0)^2+(V_0-2H_0)^2}{H_0^2},\qquad
 f_3=\frac{(U_0-H_0)^2+(V_0-2H_0)^2}{H_0^2}.
\]
Adjoin both square roots to the conic function field.  On its affine
chart the return is
\[
 x=\frac{U_0+1}{2},\quad y=\frac{V_0}{2},\quad
 (r_0,r_1,r_2,r_3)
 =\left(|R_0|,|R_1|,\frac{|\sqrt{f_2}|}{2},
                           \frac{|\sqrt{f_3}|}{2}\right).
\]

\begin{proposition}\label{prop:ell-adjacent-genus}
For every rational $c_0>0$ with $c_0\ne1$, this normalized cover
is geometrically connected of genus five.
\end{proposition}
\begin{proof}
Over $\overline\Q$, the numerator of $f_2$ factors into the two lines
$U_0+H_0\pm i(V_0-2H_0)=0$.
For the line with sign $\delta=\pm1$, evaluation of the dual conic is
\[
 \frac{1}{\lambda-1}-\frac1\lambda
 -\frac{(1-2\delta i)^2}{\lambda(\lambda-1)}
 =\frac{4+4\delta i}{\lambda(\lambda-1)}\ne0.
\]
Both intersections are therefore transverse.  The two lines meet at
$(-1,2,1)$, which would belong to the conic only if
$\lambda^2-6\lambda+1=0$; this is impossible for rational $\lambda$.
Thus $f_2$ has four distinct simple zeros.  The same calculation
with $U_0\mapsto-U_0$ proves this for $f_3$.

The two zero sets are disjoint.  Their numerator difference is
$4U_0H_0$.  At $H_0=0$ a common zero would satisfy
$U_0^2+V_0^2=0$, which together with the conic forces
$U_0=V_0=0$.  In the chart $H_0=1$, a common zero has
$U_0=0$ and $V_0=2\pm i$, whence the conic requires
$\lambda=4\pm4i$, again impossible.
The two points of the conic at infinity are distinct and neither
numerator vanishes there; both poles of each $f_i$ have order two.
The radicands are consequently independent square classes, and
the degree-four cover has eight branch values, each contributing
two to the ramification divisor.  Hence
$2g-2=4(-2)+8\cdot2=8$.
\end{proof}

\begin{theorem}[Fixed-pair finiteness]\label{thm:ell-fixed-pair-finite}
Fix two distinct vertices, signs $\varepsilon_i,\varepsilon_j\in\{1,-1\}$,
and a nonzero rational number $c_0$.
There are only finitely many strict rational four-distance points satisfying
$\varepsilon_i r_i+\varepsilon_j r_j=c_0$.
\end{theorem}
\begin{proof}
Changing both signs and applying a square symmetry reduces the pair
to one of the two constructions above.  Their return formulas show
that every strict point is the image of a rational point on the
corresponding smooth projective genus-five curve.  Faltings'
theorem \cite{faltings1983} makes this rational-point set finite.
\end{proof}

\subsection{The fibre with opposite difference one fifth}
\label{subsec:ell-one-fifth}

For $c_0=1/5$, take $\mu=1$, $\nu=5$, $\varepsilon=-1$ and
$Z=7z$ in \eqref{eq:ell-fixed-opposite}.  The curve becomes
\begin{equation}\label{eq:ell-one-fifth}
 C:\quad P^2=z^2+50,\qquad
 Q^2=z^2+14z+51,\qquad R^2=z^2-14z+51.
\end{equation}
Its strict positive branch has $P,Q,R>0$ and
\begin{equation}\label{eq:ell-one-fifth-return}
 \begin{aligned}
 x&=\tfrac12-\frac{P+7z}{100},&
 y&=\tfrac12-\frac{P-7z}{100},\\
 (r_0,r_1,r_2,r_3)&=\tfrac1{10}(P-1,Q,R,P+1),&
 |P+7z|&<50,&|P-7z|&<50.
 \end{aligned}
\end{equation}
The projective model in coordinates $[P:Q:R:z:T]$ is obtained by
homogenizing the three equations.  Above $T=0$ it has the eight
rational points $[\pm1:\pm1:\pm1:1:0]$.  These points are smooth:
in the chart $z=1$ the derivatives in $P,Q,R$ form an invertible
diagonal matrix, and $T$ is a local parameter.

\begin{proposition}[Elliptic decomposition of the fixed fibre]
\label{prop:ell-one-fifth-jacobian}
There is a $\Q$-isogeny
\[
 \Jac(C)\sim_{\Q} E_1\times F\times E_3\times E_4^2,
\]
where
\begin{align*}
 E_1&:\ y^2=x^3-4606x^2+6245001x,\\
 F&:\ y^2=x^3-194x^2+9801x,\\
 E_3&:\ y^2=x^3-9902x^2+25492401x,\\
 E_4&:\ y^2=x^3-202x^2+9801x.
\end{align*}
\end{proposition}
\begin{proof}
Put $f_0=z^2+50$, $f_\pm=z^2\pm14z+51$.
The seven nontrivial characters of the sign-change group
$(\Z/2\Z)^3$ give quotients $w^2=\prod_{j\in I}f_j$.
Their genera are $0,0,0,1,1,1,2$, according as $|I|$ is one,
two, or three.  The pullbacks of their differential spaces are
the respective character spaces in $H^0(C,\Omega_C^1)$, and
their dimensions sum to five.  The product of the associated
Jacobian pullback maps is therefore an isogeny over $\Q$.

The quotient $w^2=f_+f_-=z^4-94z^2+2601$ maps with degree two
to $y^2=x^3-94x^2+2601x$ by $(x,y)=(z^2,zw)$.
The scaling $(x,y)\mapsto(49x,343y)$ gives $E_1$.
The two quotients $w^2=f_0f_\pm$ are isomorphic by $z\mapsto-z$.
For the plus sign the substitutions
\[
 W_0=w+z^2+7z+26,\qquad
 x=2W_0+49,\qquad y=(4z+14)W_0+336
\]
give $E_4$.  Their rational inverse on $W_0\ne0$ is
\[
 W_0=(x-49)/2,\qquad
 z=\frac{y-14W_0-336}{4W_0},\qquad
 w=W_0-z^2-7z-26.
\]

The remaining genus-two curve is
\[
 D:\quad w^2=(z^2+50)(z^4-94z^2+2601).
\]
Its involutions $(z,w)\mapsto(-z,w)$ and
$(z,w)\mapsto(-z,-w)$ have elliptic quotients.
The first is $(x,y)=(z^2+50,w)$ on $F$.
For the second put $\zeta=z^2$ and $v=zw$, so that
$v^2=\zeta(\zeta+50)(\zeta^2-94\zeta+2601)$.
More generally, for
$v^2=\zeta(\zeta+c_1)(\zeta^2+a_1\zeta+b_1)$, the substitutions
\[
 X=\frac{c_1b_1}{\zeta}+b_1,\qquad
 Y=\frac{c_1b_1v}{\zeta^2}
\]
give
\[
 Y^2=X^3+(c_1a_1-2b_1)X^2
          +b_1(b_1-c_1a_1+c_1^2)X.
\]
Taking $(c_1,a_1,b_1)=(50,-94,2601)$ yields $E_3$.
The two one-dimensional differential eigenspaces on $D$ prove
$\Jac(D)\sim_\Q F\times E_3$, completing the decomposition.
\end{proof}

\begin{example}\label{ex:ell-nonlift}
The elliptic quotient $F$ contains
\[
 (\xi,\eta)=\left(
 \frac{1476225}{18496},-\frac{589188735}{2515456}\right).
\]
This follows by direct substitution in its cubic equation.
Its first coordinate is $(1215/136)^2$ and satisfies $50<\xi<99$.
A lift to \eqref{eq:ell-one-fifth} would require
\[
 Z^2=49(\xi-50)=\left(\frac{35}{136}\right)^2 22057.
\]
Since $148^2<22057<149^2$, the required rational root does not exist.
Thus the interval \eqref{eq:ell-quotient-interval} and a square
elliptic coordinate do not characterize the image of $C(\Q)$;
the three simultaneous equations in \eqref{eq:ell-one-fifth}
give the complete lift.
\end{example}
